\section{Вибір криптографічної бібліотеки для різних IT-систем}

\subsection{Сценарій А. Backend-сервіс}

\begin{itemize}
    \item Надійність та перевіреність. Сервер працює з конфеденційними даними - тому це обов'язково.
    \item Підтримка сучасних криптографічних алгоритів. Необіхідність, коли мова йде про надійні алгоритми шифрування, гешування та роботи з ключами.
    \item Якісний генератор випадкових чисел. Криптографічна стійка випадковість необхідна для генерації ключів, salt, nonce.
    \item Захист від типових помилок або атак. Ризик повинен бути зменшений.
    \item Оновлення безпеки. Попри гарну схему, важливо регулярне оновлення, щоб швидко виправити виявлені вразливості.
\end{itemize}

\subsection{Сценарій Б. Мобільний застосунок}

\begin{itemize}
    \item Надійне шиврування локальних даних. Якщо зловмисник отрумає доступ до файлів застосунку, то конфіденційна інформація повинна бути збережена.
    \item Низьке споживання пам'яті. Мобільні пристрої більш обмежені, аніж сервери.
    \item Енергоефективність. Надмірна кількість операцій може швидко зменшувати заряд мобільного пристрою.
    \item Безпечне очищення пам'яті. Після використання секретні дані бажано видалити з пам'яті.
    \item Безпечне зберігання ключів. Ключ не можна просто залишити у файлі або в коді застосунку.
\end{itemize}

\subsection{Сценарій В. Сервіс цифрового підпису}

\begin{itemize}
    \item Підтримка стандарту PKCS\#11 та взаємодії з HSM (Hardware Security Module). Оскільки закритий ключ є найціннішим активом, він ніколи не повинен потрапляти у відкритому вигляді в оперативну пам'ять сервера; бібліотека має вміти делегувати операцію підпису зовнішньому апаратному модулю безпеки.
    \item Захист від атак побічними каналами (Constant-Time Implementation). Реалізація алгоритмів підпису (RSA, ECDSA, Ed25519) повинна мати постійний час виконання та використовувати техніку криптографічного засліплення (\textit{blinding}) для виключення витоку частин закритого ключа через вимірювання часу чи енергоспоживання.
    \item Стійкість до дефектів обчислень (Fault Injection Resistance). Під час використання оптимізації CRT для RSA бібліотека обов'язково повинна верифікувати підпис перед відправленням клієнту, запобігаючи витоку множників ключа через обчислювальні збої.
    \item Сувора відповідність нормативним стандартам і FIPS-сертифікація. Сервіси цифрового підпису зазвичай мають юридичну силу, тому криптографічний модуль повинен мати сертифікат відповідності стандартам (FIPS 140-3, eIDAS, Common Criteria).
\end{itemize}

Обраний засіб: Програмний інтерфейс на базі PKCS\#11 (наприклад, \textit{OpenSSL} із відповідним engine/provider або \textit{SunPKCS11}) у зв'язці з апаратним модулем безпеки HSM.

\subsection{Сценарій Г. Embedded/IoT-пристрій}

\begin{itemize}
    \item Мінімальний розмір коду (Small Footprint) та відсутність динамічного виділення пам'яті. Бібліотека повинна займати мінімальний обсяг ROM/Flash і RAM, дозволяти статичне виділення буферів та уникати фрагментації обмеженої динамічної пам'яті (heap).
    \item Підтримка апаратних криптографічних блоків мікроконтролера. Можливість задіяти вбудовані в кристал ESP32 апаратні прискорювачі AES, SHA та RSA/ECC для прискорення обчислень та збереження заряду акумулятора.
    \item Використання алгоритмів на еліптичних кривих (ECC). Застосування Ed25519 або ECDSA замість ресурсомісткого RSA, оскільки еліптична криптографія забезпечує високу стійкість за значно меншої довжини ключів (256 біт замість 2048 біт) і менших витрат оперативної пам'яті.
    \item Підтримка апаратного джерела ентропії (TRNG). Для безпечної генерації ключів та nonce бібліотека повинна отримувати випадкові значення безпосередньо з апаратного генератора шуму ESP32, а не використовувати звичайні некриптографічні PRNG.
\end{itemize}

Обраний засіб: Бібліотека \textit{Mbed TLS}, інтегрована до складу \textit{ESP-IDF}, з увімкненою апаратною підтримкою мікроконтролера ESP32.
